\clearpage\thispagestyle{empty}
\pdfbookmark[0]{Quellen zur Besprechung von Beer und Mädler}{beermaedler1841-sources}
\begin{center}\Large\bfseries Quellen zur Besprechung von Beer und Mädler\end{center}

Die Übertragung folgt Carl Friedrich Gauß’ Besprechung in den \emph{Werken}, Bd.~VI (Göttingen, 1874), S.~652--654. Der Erstabdruck erschien in den \emph{Göttingischen gelehrten Anzeigen}, Stück~137, 30.~August 1841, S.~1361--1365. Besprochen werden Beer und Mädlers \emph{Beiträge zur physischen Kenntniss der himmlischen Körper im Sonnensysteme} (Weimar, B. F. Voigt, 1841) und ihre \emph{Fragments sur les corps célestes du système solaire} (Paris, Bachelier, 1840).

KI-Übertragung aus dem historischen Deutsch in modernes Deutsch. KI-erstellter Satz.

\href{https://gdz.sub.uni-goettingen.de/id/PPN236007467}{Werkausgabe}. Erstdruck, S.~\href{https://images.p.sub.uni-goettingen.de/iiif/image/gdz:PPN385030444_1841_2:00000355/full/full/0/default.jpg}{1361}, \href{https://images.p.sub.uni-goettingen.de/iiif/image/gdz:PPN385030444_1841_2:00000356/full/full/0/default.jpg}{1362}, \href{https://images.p.sub.uni-goettingen.de/iiif/image/gdz:PPN385030444_1841_2:00000357/full/full/0/default.jpg}{1363}, \href{https://images.p.sub.uni-goettingen.de/iiif/image/gdz:PPN385030444_1841_2:00000358/full/full/0/default.jpg}{1364}, \href{https://images.p.sub.uni-goettingen.de/iiif/image/gdz:PPN385030444_1841_2:00000359/full/full/0/default.jpg}{1365}. \href{https://doi.org/10.3931/e-rara-49144}{Weimarer Ausgabe von 1841}. \href{https://books.google.nl/books?id=F85CAQAAMAAJ&pg=PP7}{Pariser Ausgabe von 1840}.

\clearpage
